\documentclass[10pt]{article}

%%% These are some packages that are useful
\usepackage{amsmath,amssymb, amscd,amsbsy, amsthm, enumerate}
\usepackage[export]{adjustbox}
\usepackage{lastpage}
\usepackage[top=1in, bottom=1in, left=1in, right=1in]{geometry}
\usepackage[unicode]{hyperref}
\usepackage{tikz, pgfplots, xcolor, fancyhdr}
\usepackage{multicol,caption}
\usepackage{lipsum}
\usepackage[version=4]{mhchem}
\usepackage{float}
\usepackage{enumitem}
\usepackage{xurl}

%%% Page formatting
%\setlength{\headsep}{30pt}
\setlength{\textheight}{9in}
\newcommand{\tab}{\hspace{1cm}}
%\setlength{\parindent}{25pt}

\title{Know Thy Enemy (Part 7) 	: Hidden In Plain Sight}
\author{Antonius Torode}

%%% Header and Footer Info
\pagestyle{fancy}
\fancyhead[L]{{\large Template - \textbf{Change 03}}}
\fancyhead[C]{\today}
\fancyhead[R]{Name: Antonius Torode}


\fancyhf{} % sets both header and footer to nothing
\renewcommand{\headrulewidth}{0pt}
% your new footer definitions here

\fancyfoot[L]{}
\fancyfoot[C]{}
\fancyfoot[R]{\thepage\ of \pageref{LastPage}}

% Used to define spacing and format of References
\let\OLDthebibliography\thebibliography
\renewcommand\thebibliography[1]{
	\OLDthebibliography{#1}
	\setlength{\parskip}{0pt}
	\setlength{\itemsep}{0pt plus 0.3ex}
}

\newenvironment{Figure}
  {\par\medskip\noindent\minipage{\linewidth}}
  {\endminipage\par\medskip}
  
\renewcommand{\abstractname}{Introduction}

%%% Document Starts now
\begin{document}

\maketitle
\thispagestyle{fancy}


\begin{abstract}
This message explores the close ties between the \textit{enemy} and symbolism. Throughout history (heavily including modern times), the depiction of \textit{demons} has been closely related to symbolism. This message is written in such a way to be delivered in conjunction with a collection of images. These images are compiled in an associated power-point presentation which will be referenced throughout. A power-point link can be found as the first listed reference \cite{ppt}.
\end{abstract}

\begin{multicols}{2}

One common things that comes to mind when people hear the term \textit{demon} is that of symbols; whether it's spell-circles, pentagrams, the Baphomet symbol, a golden calf, or something else. One reason this is the case is because many forms of entertainment use these concepts in conjunction. Whether it's a scary movie, a Marvel movie, or even a video game. The depiction of \textit{demons} is very commonly found closely tied to the depiction of symbols. 

\begin{thebibliography}{9}
{\footnotesize
	\bibitem{ppt} TODO
}
\end{thebibliography}

\end{multicols}


%%%%%%%%%%%%%%%%%%%%%%%%%%%%%%%%%%%%%%%%%%%%%%%%%%%%%%%%%%%%%%%%%%%%%%%%%%%%%%%%%%%%%%%%%%%
\end{document}





















}{den}